\section{与审稿员预期批评的预设回答}

以下针对审稿人可能提出的六个典型问题逐一给出回应。
问题1：Meta-label 会显著增加工作量吗？
回答：不会。Scope 是必要决策字段；Difficulty 与 \texttt{model\_issue} 属于修改时的同步记录。额外时间预计 $<10$ 秒，占平均标注时长约 8\%。
问题2：多选标签如何定义共识？
回答：采用 Per-tag 二元共识。每个 tag 独立多数票，报告 per-tag Fleiss' Kappa。为避免 ``空值=无困难'' 与 ``漏填'' 不可区分，采集协议在提交时对多选元标签执行硬校验：至少选 1 项（空选不可提交）且禁止两类互斥共选（\texttt{trivial}+其他困难标签；\texttt{acceptable}+其他失败模式）。NA 仍可能因系统/导出缺损出现，因此 NA 不计入有效人数并单独披露缺失率；同时披露被拦截/拒收次数与原因分布作为过程证据。
问题3：样本量有限，结论是否可靠？
回答：主终点预注册，使用 BCa bootstrap CI，报告效应量与 MDE，不以 $p$ 值作为唯一依据；同时披露各场景样本量与 CI 宽度。
问题4：离线路由模拟是否弱于在线实验？
回答：离线路由是当前阶段可复现且成本可控的标准做法，可用于形成在线实验先验；本文明确其边界并保留后续在线验证。
问题5：加权共识是否循环论证，OOD 下专家会否被淹没？
回答：权重锚点来自 PreScreen 测试样本与专家参考标准，并采用 leave-one-task-out，避免同任务自证；$w_{max}$ 在 PreScreen 阶段按预注册的方案 A（不足则方案 B）锁定后冻结，避免事后调参；高 $d_{t}$ 场景进入 stress mode，提高冗余或触发专家复核。
问题6：改进幅度是否过小？
回答：以预注册 MDE 判定工程意义；只有效应量跨过 MDE 且统计证据一致时，才声称改进成立。


